\section{Системы с фазовым обслуживанием. Модель $M|Er_k|1$}

\begin{definition}{Система $M|Er_k|1$}
Одноканальная система массового обслуживания, на вход которой поступает пуассоновский поток заявок с интенсивностью $\lambda$. Процесс обслуживания каждой заявки состоит из $k$ последовательных независимых этапов (фаз). Время прохождения каждой фазы распределено по экспоненциальному закону. Суммарное время обслуживания заявки $S$ подчиняется распределению Эрланга $k$-го порядка ($S \sim Er_k$).
\end{definition}

Для обеспечения математического ожидания полного времени обслуживания $\E[S] = \frac{1}{\mu}$, среднее время выполнения одной обособленной фазы должно составлять $\frac{1}{k\mu}$. Следовательно, интенсивность переходов процесса обслуживания между смежными фазами равна $k\mu$.

\begin{figure}[h!]
\begin{center}
\begin{tikzpicture}[>=Stealth, scale=1.0, every node/.style={font=\small}]
    % Определение кастомных стилей для фазовых переходов (толщина 1.0pt)
    \tikzset{redarrow/.style={->, dashed, red!85!black, line width=1pt}}
    \tikzset{rednode/.style={circle, draw, dashed, red!85!black, line width=1.5pt, minimum size=0.6cm, inner sep=1pt}}

    % Состояние 0 (придвинуто к красной стрелке, которая заканчивается на 0.9)
    \node[circle, draw, thick, minimum size=0.8cm] (S0) at (0.5,0) {0};

    % Макро-состояние 1 (Смещено вправо, левая граница на x=2.0)
    \node[draw, thick, rounded corners=4pt, minimum width=4.8cm, minimum height=1.6cm] (B1) at (4.4, 0) {};
    \node[circle, draw, thick, minimum size=0.8cm] (M1) at (4.4, -1.8) {1};

    % Макро-состояние 2
    \node[draw, thick, rounded corners=4pt, minimum width=4.8cm, minimum height=1.6cm] (B2) at (10.8, 0) {};
    \node[circle, draw, thick, minimum size=0.8cm] (M2) at (10.8, -1.8) {2};

    % Многоточие уходящее вправо
    \node at (15.2, 0) {\Large $\dots$};

    % Фазы в состоянии 1
    \node[rednode] (P1) at (2.6,0) {1};
    \node[rednode] (P2) at (3.8,0) {2};
    \node[red!85!black] (Pdots1) at (5.0,0) {\dots};
    \node[rednode] (Pk) at (6.2,0) {$k$};

    % Фазы в состоянии 2
    \node[rednode] (Pk1) at (9.0,0) {\scriptsize $k\!\!+\!\!1$};
    \node[rednode] (Pk2) at (10.2,0) {\scriptsize $k\!\!+\!\!2$};
    \node[red!85!black] (Pdots2) at (11.4,0) {\dots};
    \node[rednode] (P2k) at (12.6,0) {\scriptsize $2k$};

    % Внутренние красные пунктирные стрелки
    \draw[redarrow] (P2) -- node[above, font=\scriptsize] {$k\mu$} (P1);
    \draw[redarrow] (Pdots1) -- node[above, font=\scriptsize] {$k\mu$} (P2);
    \draw[redarrow] (Pk) -- node[above, font=\scriptsize] {$k\mu$} (Pdots1);

    \draw[redarrow] (Pk1) -- node[above, font=\scriptsize] {$k\mu$} (Pk);
    \draw[redarrow] (Pk2) -- node[above, font=\scriptsize] {$k\mu$} (Pk1);
    \draw[redarrow] (Pdots2) -- node[above, font=\scriptsize] {$k\mu$} (Pk2);
    \draw[redarrow] (P2k) -- node[above, font=\scriptsize] {$k\mu$} (Pdots2);

    % Внешние пунктирные стрелки
    \draw[redarrow] (P1.west) -- node[above, font=\scriptsize] {$k\mu$} (0.9, 0);
    \draw[redarrow] (14.3, 0) -- node[above, font=\scriptsize] {$k\mu$} (P2k.east);

    % Макроскопические черные стрелки (\lambda)
    % Первая стрелка теперь "висит в воздухе" симметрично нижней
    \draw[->, thick] (0.9, 0.8) to [bend left=25] node[above] {$\lambda$} (B1.north);
    \draw[->, thick] (B1.north) to [bend left=25] node[above] {$\lambda$} (B2.north);
    \draw[->, thick] (B2.north) to [bend left=25] node[above] {$\lambda$} (14.3, 0.8);

    % Макроскопические черные стрелки (\mu)
    \draw[->, thick] (14.3, -0.8) to [bend left=25] node[below] {$\mu$} (B2.south);
    \draw[->, thick] (B2.south) to [bend left=25] node[below] {$\mu$} (B1.south);
    \draw[->, thick] (B1.south) to [bend left=25] node[below] {$\mu$} (0.9, -0.8);
\end{tikzpicture}
\caption{Структура фазового обслуживания системы $M|Er_k|1$. Состояния марковской цепи сгруппированы в макро-состояния (соответствующие числу заявок в системе), внутри которых протекает последовательное выполнение независимых фаз. Нумерация микро-состояний ведется по суммарному числу оставшихся этапов обслуживания.}
\end{center}
\end{figure}

Состояние системы в произвольный момент времени описывается двумерным вектором $(n, i)$, где $n$ — общее число заявок в системе, а $i \in \{1, 2, \dots, k\}$ — номер текущей фазы обслуживаемой заявки. 

Для снижения размерности задачи сведем двумерное пространство состояний к одномерному. Введем случайную величину $J$, равную суммарному числу невыполненных фаз обслуживания для всех заявок, находящихся в системе. При наличии $n$ заявок и обслуживании первой из них на стадии $i$, искомое значение $J$ задается соотношением:
$$ J = (n - 1)k + i, \quad \text{при } n > 0, $$
и $J = 0$ при отсутствии заявок в системе ($n = 0$).

\begin{statement}
Марковская цепь, описывающая изменение суммарного числа невыполненных фаз $J$ в системе $M|Er_k|1$, стохастически эквивалентна марковской цепи системы с групповым поступлением заявок $M^{[X]}|M|1$, в которой размер группы строго фиксирован и равен $k$ ($X = k = \text{const}$).
\end{statement}

\begin{sproof}
Рассмотрим динамику изменения состояния системы, выраженного через переменную $J(t)$. Поступление новой заявки образует пуассоновский поток с интенсивностью $\lambda$ и мгновенно увеличивает суммарное число невыполненных фаз на $k$. Данное событие соответствует скачкообразному переходу марковской цепи из состояния $j$ в $j + k$.

В свою очередь, обслуживание заявок представляет собой последовательное выполнение фаз. Завершение обработки каждой отдельной фазы происходит с интенсивностью $k\mu$ и строго уменьшает текущее число оставшихся фаз на единицу (переход $j \to j - 1$). Таким образом, процесс изменения числа фаз стохастически идентичен процессу изменения числа элементов в системе $M^{[X=k]}|M|1$, где пакеты фиксированного размера $k$ поступают с интенсивностью $\lambda$, а элементы пакета обслуживаются поодиночке с интенсивностью $k\mu$.
\end{sproof}

\begin{figure}[h!]
\begin{center}
\begin{tikzpicture}[
    scale=1.5,
    >={Stealth[length=3mm]},
    state/.style={circle, draw, minimum size=0.8cm, thick},
    every node/.style={font=\small}
]
    % Состояния
    \node[state] (s0) at (0,0) {0};
    \node[state] (s1) at (1.5,0) {1};
    \node[state] (s2) at (3,0) {2};
    \node (dots1) at (4.5,0) {\Large $\dots$};
    \node[state] (sk) at (6,0) {$k$};
    \node[state] (sk1) at (7.5,0) {\scriptsize $k\!+\!1$};
    \node (dots2) at (9,0) {\Large $\dots$};

    % Обслуживание (k\mu) - идет влево (уменьшение фаз)
    \draw[->, thick] (s1) to [bend left=30] node[midway, below] {$k\mu$} (s0);
    \draw[->, thick] (s2) to [bend left=30] node[midway, below] {$k\mu$} (s1);
    \draw[->, thick] (dots1) to [bend left=30] node[midway, below] {$k\mu$} (s2);
    \draw[->, thick] (sk) to [bend left=30] node[midway, below] {$k\mu$} (dots1);
    \draw[->, thick] (sk1) to [bend left=30] node[midway, below] {$k\mu$} (sk);

    % Прибытие (\lambda) - прыжок на k вправо (приход k фаз)
    \draw[->, thick] (s0) to [bend left=45] node[midway, above] {$\lambda$} (sk);
    \draw[->, thick] (s1) to [bend left=45] node[midway, above] {$\lambda$} (sk1);
    \draw[->, thick] (s2) to [bend left=45] node[midway, above] {$\lambda$} (dots2);
\end{tikzpicture}
\caption{Граф интенсивностей переходов для суммарного числа невыполненных фаз $J$ в системе $M|Er_k|1$. Граф полностью эквивалентен модели $M^{[X=k]}|M|1$ с индивидуальной интенсивностью обслуживания $k\mu$.}
\end{center}
\end{figure}

Ключевое отличие $M|Er_k|1$ от классической системы с групповым поступлением заключается в физической интерпретации состояний: в $M^{[X=k]}|M|1$ окончание одной стадии обслуживания означает уход заявки из системы, тогда как в $M|Er_k|1$ заявка покидает прибор только после завершения всех $k$ стадий. Данная связь позволяет напрямую вывести временные характеристики системы.

\begin{theorem}{Макроскопические характеристики $M|Er_k|1$}
При выполнении условия стационарности $\rho = \frac{\lambda}{\mu} < 1$, метрики эффективности системы определяются следующими соотношениями:
\begin{itemize}
    \item Среднее время ожидания заявки в очереди: $\displaystyle W_q = \frac{k+1}{2k} \cdot \frac{\rho}{\mu(1 - \rho)}$.
    \item Среднее число заявок в очереди: $\displaystyle L_q = \frac{k+1}{2k} \cdot \frac{\rho^2}{1 - \rho}$.
    \item Среднее время пребывания заявки в системе: $\displaystyle W = W_q + \frac{1}{\mu}$.
    \item Среднее число заявок в системе: $\displaystyle L = L_q + \rho$.
\end{itemize}
\end{theorem}

\begin{tproof}
Вычислим среднее время ожидания заявки в очереди $W_q$. При использовании дисциплины обслуживания FCFS (First-Come, First-Served) время ожидания вновь прибывшей заявки строго равно суммарному времени, необходимому для завершения всех фаз обслуживания, уже находящихся в системе к моменту ее поступления.

В силу свойства PASTA (Poisson Arrivals See Time Averages), пуассоновский поток заявок застает систему в стационарном состоянии. Следовательно, математическое ожидание числа невыполненных фаз, наблюдаемое прибывающей заявкой, совпадает со средним числом фаз в произвольный момент времени $\E[J]$.

Для вычисления $\E[J]$ воспользуемся установленной стохастической эквивалентностью с системой $M^{[X]}|M|1$. Подставим параметры эквивалентной модели в формулу для средней длины очереди системы с групповым поступлением:
\begin{itemize}
    \item Интенсивность входного потока групп: $\lambda$.
    \item Размер группы: $X = k$ (детерминированная величина, $\E[X] = k$, $\E[X^2] = k^2$).
    \item Интенсивность обслуживания одного элемента (одной фазы): $\mu_{\text{phase}} = k\mu$.
    \item Коэффициент загрузки по фазам: $\rho_{\text{phase}} = \frac{\lambda \E[X]}{\mu_{\text{phase}}} = \frac{\lambda k}{k\mu} = \rho$.
\end{itemize}

Среднее число фаз в системе равно:
$$ \E[J] = \frac{\rho_{\text{phase}} + \frac{\lambda}{\mu_{\text{phase}}} \E[X^2]}{2(1 - \rho_{\text{phase}})} = \frac{\rho + \frac{\lambda}{k\mu} k^2}{2(1 - \rho)} = \frac{\rho + \rho k}{2(1 - \rho)} = \frac{k+1}{2} \cdot \frac{\rho}{1 - \rho}. $$

Поскольку среднее время выполнения одной фазы составляет $\frac{1}{k\mu}$, полное время ожидания вычисляется как произведение среднего числа фаз на длительность одной фазы:
$$ W_q = \E[J] \cdot \frac{1}{k\mu} = \frac{k+1}{2} \cdot \frac{\rho}{1 - \rho} \cdot \frac{1}{k\mu} = \frac{k+1}{2k} \cdot \frac{\rho}{\mu(1 - \rho)}. $$

Дальнейший вывод метрик базируется на фундаментальных законах теории массового обслуживания. Применяя закон Литтла для подсистемы очереди ($L_q = \lambda W_q$), получаем:
$$ L_q = \lambda \cdot \frac{k+1}{2k} \cdot \frac{\rho}{\mu(1 - \rho)} = \frac{k+1}{2k} \cdot \frac{\rho^2}{1 - \rho}. $$

Среднее время пребывания заявки в системе складывается из времени ожидания и среднего времени обслуживания:
$$ W = W_q + \E[S] = W_q + \frac{1}{\mu}. $$

Среднее число заявок в системе в целом определяется по закону Литтла ($L = \lambda W$):
$$ L = \lambda \left( W_q + \frac{1}{\mu} \right) = L_q + \rho. $$
\end{tproof}

\begin{note}
Коэффициент $\frac{k+1}{2k}$ в полученных формулах можно представить в виде $\frac{1}{2}\left(1 + \frac{1}{k}\right)$. Рассмотрим асимптотическое поведение системы:
\begin{itemize}
    \item При $k = 1$ модель вырождается в классическую $M|M|1$, а множитель обращается в 1.
    \item При $k \to \infty$ распределение Эрланга стремится к детерминированной величине (распределение Дирака). Модель переходит в $M|D|1$, а множитель асимптотически стремится к $\frac{1}{2}$. 
\end{itemize}
Данный предел аналитически согласуется со следствиями из формулы Полячека-Хинчина, согласно которым полное устранение стохастичности времени обслуживания приводит к двукратному сокращению длины очереди при фиксированном коэффициенте загрузки.
\end{note}

\section{Протоколы множественного доступа. Алгоритмы ALOHA}

Рассмотрим сетевую архитектуру, использующую два независимых частотных канала: прямую линию связи (Downlink) и обратную линию связи (Uplink). В канале Downlink центральный узел осуществляет широковещательную передачу пакетов. В канале Uplink множество из $n$ независимых терминалов конкурируют за доступ к разделяемой среде передачи.

В классическом протоколе ALOHA терминал инициирует передачу в канал немедленно после генерации пакета, без предварительного контроля состояния среды.

\begin{note}
Отсутствие механизма прослушивания несущей (Carrier Sense) в классической архитектуре ALOHA обусловлено требованием максимальной аппаратной простоты терминалов. На этапе разработки протокола генерируемый трафик имел низкую общую интенсивность, и вероятностью коллизий пренебрегали. Механизмы предварительного контроля состояния канала (семейство CSMA) были формализованы позднее как ответ на рост нагрузки в сетях.
\end{note}

Поскольку передачи узлов независимы, пакеты могут пересекаться во времени. Наложение сигналов называется коллизией, в результате которой информация искажается и не может быть демодулирована. Для обеспечения гарантированной доставки применяется механизм автоматического запроса повторения (ARQ): при успешном приеме центральный узел отправляет кадр подтверждения (ACK). Отсутствие подтверждения диагностируется как коллизия и инициирует повторную передачу пакета через случайный интервал времени.

\subsection{Асинхронная ALOHA}

Определим абсолютную пропускную способность канала $S_{\text{abs}}$ как среднее число успешно переданных пакетов в единицу времени. Зафиксируем параметры модели:

\begin{itemize}
    \item В зоне радиовидимости базовой станции находятся $n$ передатчиков. Мощности принимаемых сигналов идентичны (эффект захвата отсутствует).
    \item Все генерируемые пакеты имеют строго фиксированную длину. Время передачи одного пакета составляет $\Delta t$.
    \item Суммарный поток пакетов (включая первичные и повторные передачи), генерируемый всеми узлами, является пуассоновским с абсолютной интенсивностью $\lambda$ [пакетов/с].
    \item Емкость буфера каждого терминала равна $1$. При наличии пакета в буфере новые заявки отбрасываются.
    \item Канал является идеальным: единственной причиной потери пакета выступает коллизия.
\end{itemize}

Введем безразмерную величину $G \dot= \lambda \Delta t$, определяющую предложенную нагрузку на канал (среднее число пакетов, поступающих за время $\Delta t$). При симметричной нагрузке интенсивность трафика отдельного узла равна $\frac{\lambda}{n}$. Рассмотрим асимптотический режим, при котором $n \to \infty$, а суммарная интенсивность $\lambda$ конечна и фиксирована.

\begin{statement}[Пропускная способность чистой ALOHA при бесконечном числе устройств]
В асимптотическом режиме $n \to \infty$ нормированная пропускная способность $S = S_{\text{abs}} \Delta t$ асинхронной сети ALOHA определяется выражением:
$$S = G e^{-2G}$$
Глобальный максимум достигается при $G = \frac{1}{2}$ и составляет $S_{\max} = \frac{1}{2e} \approx 0.184$.
\end{statement}

\begin{sproof}
Абсолютная пропускная способность выражается через интенсивность входного потока и вероятность отсутствия коллизии: $S_{\text{abs}} = \lambda \cdot \mathbb{P}(\text{успех})$. 

Для вычисления вероятности успешной передачи зафиксируем некоторый целевой пакет, инициированный в произвольный момент времени $t_0$. Поскольку длительность передачи равна $\Delta t$, она завершится в момент $t_0 + \Delta t$. Передача окажется успешной тогда и только тогда, когда ни один другой узел не начнет передачу в интервале $[t_0 - \Delta t, t_0 + \Delta t]$. Длина этого «уязвимого интервала» (vulnerable period) строго равна $\tau = 2\Delta t$.

\begin{center}
\begin{tikzpicture}[
    scale=1.1,
    >={Stealth[length=2.5mm]},
    every node/.style={font=\small},
    packet/.style={draw=black, thick, minimum height=0.8cm, align=center}
]
    % --- ФОНОВАЯ ПОДСВЕТКА ОПАСНОГО ИНТЕРВАЛА ---
    \fill[BG] (2.5, -0.8) rectangle (7.5, 6.2);

    % --- ОСИ И ПОДПИСИ СТАНЦИЙ ---
    % Оси начинаются левее (-0.5), чтобы вместить успешный пакет
    \foreach \y/\label in {4.5/1, 3.0/2, 1.5/3, 0.0/4} {
        \node[left, font=\bfseries] at (-0.8, \y) {Станция \label};
        \draw[->, thick] (-0.5, \y) -- (11.5, \y) node[right] {$t$};
    }

    % --- ВЕРТИКАЛЬНЫЕ ЛИНИИ (ГРАНИЦЫ ИНТЕРВАЛОВ) ---
    \draw[thick, dashed, DeepSkyBlue4] (2.5, -0.5) -- (2.5, 6.2) node[above, black] {$t_0 - \Delta t$};
    \draw[thick, dashed, black] (5.0, -0.5) -- (5.0, 6.2) node[above, black] {$t_0$};
    \draw[thick, dashed, DeepSkyBlue4] (7.5, -0.5) -- (7.5, 6.2) node[above, black] {$t_0 + \Delta t$};

    % --- МАРКЕР УЯЗВИМОГО ВРЕМЕНИ НАД ГРАФИКОМ ---
    \draw[<->, thick] (2.5, 5.8) -- (7.5, 5.8) node[midway, fill=BG, inner sep=2pt, font=\footnotesize\bfseries] {Уязвимый интервал ($2\Delta t$)};

    % ==========================================
    % ОТРИСОВКА ПАКЕТОВ (\Delta t = 2.5)
    % ==========================================

    % --- СТАНЦИЯ 1 (Левая коллизия) ---
    \fill[DeepSkyBlue4!15] (3.5, 4.1) rectangle (6.0, 4.9);
    \fill[Red3!25] (5.0, 4.1) rectangle (6.0, 4.9); % Зона перекрытия
    \draw[thick] (3.5, 4.1) rectangle (6.0, 4.9);
    \node[font=\footnotesize] at (4.25, 4.5) {Пакет $i$};
    
    % Стрелка момента времени t_i
    \draw[->, thick] (3.5, 3.5) node[below, font=\scriptsize] {$t_i$} -- (3.5, 4.05);
    
    % Указатель коллизии
    \draw[<-, thick, Red3] (5.5, 4.9) -- (5.8, 5.3) node[right, font=\footnotesize\bfseries] {Коллизия};

    % --- СТАНЦИЯ 2 (Целевой пакет) ---
    \fill[SpringGreen4!15] (5.0, 2.6) rectangle (7.5, 3.4);
    \fill[Red3!25] (5.0, 2.6) rectangle (6.0, 3.4); % Перекрытие от Станции 1
    \fill[Red3!25] (6.5, 2.6) rectangle (7.5, 3.4); % Перекрытие от Станции 3
    \draw[thick] (5.0, 2.6) rectangle (7.5, 3.4);
    
    % Подпись вынесена за пределы пакета (сверху) для читаемости
    \node[font=\footnotesize\bfseries, above] at (6.25, 3.45) {Целевой пакет};
    
    % Размерная линия \Delta t
    \draw[<->, thick] (5.0, 2.3) -- (7.5, 2.3) node[midway, below, font=\footnotesize] {$\Delta t$};

    % --- СТАНЦИЯ 3 (Правая коллизия) ---
    \fill[DeepSkyBlue4!15] (6.5, 1.1) rectangle (9.0, 1.9);
    \fill[Red3!25] (6.5, 1.1) rectangle (7.5, 1.9); % Зона перекрытия
    \draw[thick] (6.5, 1.1) rectangle (9.0, 1.9);
    \node[font=\footnotesize] at (8.25, 1.5) {Пакет $j$};
    
    % Стрелка момента времени t_j
    \draw[->, thick] (6.5, 0.5) node[below, font=\scriptsize] {$t_j$} -- (6.5, 1.05);
    
    % Указатель коллизии
    \draw[<-, thick, Red3] (7.0, 1.1) -- (7.5, 0.6) node[right, font=\footnotesize\bfseries] {Коллизия};

    % --- СТАНЦИЯ 4 (Безопасная передача) ---
    % Пакет сдвинут влево и заканчивается на 2.2, чтобы не касаться границы 2.5
    \fill[SpringGreen4!15] (-0.3, -0.4) rectangle (2.2, 0.4);
    \draw[thick] (-0.3, -0.4) rectangle (2.2, 0.4);
    \node[font=\footnotesize, align=center] at (0.95, 0) {Успешная\\передача};

\end{tikzpicture}
\end{center}

В силу независимости терминалов, процесс генерации пакетов \textit{остальными} $n-1$ узлами формирует фоновый пуассоновский поток. При $n \to \infty$ интенсивность этого фонового потока совпадает с суммарной интенсивностью $\lambda$. Фиксированный целевой пакет выступает заданным условием наблюдения, поэтому он исключается из множества случайных событий фонового процесса.

Вероятность успешной передачи совпадает с вероятностью того, что на уязвимом интервале $\tau = 2\Delta t$ наступит ровно ноль событий фонового потока. Применяя формулу Пуассона для $k=0$ и среднего числа событий $\lambda \tau = \lambda(2\Delta t) = 2G$:
$$\mathbb{P}(\text{успех}) = \mathbb{P}(k=0) = \frac{(2G)^0}{0!} e^{-2G} = e^{-2G}.$$
Абсолютная пропускная способность составляет $S_{\text{abs}} = \lambda e^{-2G}$. Переходя к нормированной величине $S = S_{\text{abs}} \Delta t$ (числу успешных пакетов за время $\Delta t$), получаем итоговую формулу:
$$S = (\lambda \Delta t) e^{-2G} = G e^{-2G}.$$
Вычислим производную функции $S(G)$ для нахождения точки оптимума:
$$\frac{dS}{dG} = e^{-2G} - 2Ge^{-2G} = e^{-2G}(1 - 2G) = 0 \implies G = \frac{1}{2}.$$
\end{sproof}

\begin{note}
В дальнейших выкладках стандартной практикой является нормировка оси времени, при которой $\Delta t \equiv 1$. В этом случае понятия абсолютной и приведенной интенсивностей совпадают ($\lambda = G$, $S_{\text{abs}} = S$), а единицей измерения времени становится длительность одного пакета.
\end{note}

\subsection{Синхронная (слотированная) ALOHA при бесконечном числе устройств}

Основным фактором снижения эффективности чистой ALOHA является избыточная длина уязвимого интервала ($2\Delta t$). Алгоритм синхронной ALOHA устраняет частичные наложения пакетов за счет квантования времени.

Ось времени разбивается на равные слоты длительностью $\Delta t$. Терминалам разрешается инициировать передачу исключительно на границах слотов. Для координации узлов применяется глобальный синхронизирующий сигнал.

\begin{statement}[Пропускная способность слотированной ALOHA]
В асимптотическом режиме $n \to \infty$ нормированная пропускная способность синхронной сети ALOHA задается выражением:
$$S = G e^{-G}$$
Глобальный максимум достигается при $G = 1$ и составляет $S_{\max} = \frac{1}{e} \approx 0.368$.
\end{statement}

\begin{sproof}
Пусть передача целевого пакета запланирована в слоте, ограниченном моментами $k\Delta t$ и $(k+1)\Delta t$. Заявки, инициирующие эту передачу, формируются терминалом на предшествующем интервале. Поскольку любые передачи стартуют строго по границам слотов, частичное наложение пакетов невозможно: пакеты либо полностью совпадают во времени, либо не пересекаются вовсе.

\begin{center}
\begin{tikzpicture}[
    scale=1.1,
    >={Stealth[length=2.5mm]},
    every node/.style={font=\small},
    packet/.style={draw=black, thick, minimum height=0.8cm, align=center}
]
    % --- ФОНОВАЯ ПОДСВЕТКА СЛОТОВ ---
    % В Slotted ALOHA уязвимым является только сам слот целевого пакета
    \fill[BG] (5.0, -0.8) rectangle (7.5, 6.2);

    % --- ОСИ И ПОДПИСИ СТАНЦИЙ ---
    \foreach \y/\label in {4.5/1, 3.0/2, 1.5/3, 0.0/4} {
        \node[left, font=\bfseries] at (-0.8, \y) {Станция \label};
        \draw[->, thick] (-0.5, \y) -- (11.5, \y) node[right] {$t$};
    }

    % --- ВЕРТИКАЛЬНЫЕ ЛИНИИ (ГРАНИЦЫ СЛОТОВ) ---
    \draw[thick, dashed, black!50] (0.0, -0.5) -- (0.0, 6.2);
    \draw[thick, dashed, DeepSkyBlue4] (2.5, -0.5) -- (2.5, 6.2) node[above, black] {$t_0 - \Delta t$};
    \draw[thick, dashed, black] (5.0, -0.5) -- (5.0, 6.2) node[above, black] {$t_0$};
    \draw[thick, dashed, DeepSkyBlue4] (7.5, -0.5) -- (7.5, 6.2) node[above, black] {$t_0 + \Delta t$};
    \draw[thick, dashed, black!50] (10.0, -0.5) -- (10.0, 6.2) node[above, black] {$t_0 + 2\Delta t$};

    % --- ПОДПИСИ НОМЕРОВ СЛОТОВ ВНИЗУ ---
    \node[below, font=\footnotesize, text=gray!80!black] at (1.25, -0.6) {Слот $k-2$};
    \node[below, font=\footnotesize, text=Red3] at (3.75, -0.6) {Слот $k-1$ (Уязвимый)};
    \node[below, font=\footnotesize\bfseries, text=DeepSkyBlue4] at (6.25, -0.6) {Слот $k$};
    \node[below, font=\footnotesize, text=DeepSkyBlue4] at (8.75, -0.6) {Слот $k+1$};

    % --- МАРКЕР УЯЗВИМОГО ВРЕМЕНИ НАД ГРАФИКОМ ---
    \draw[<->, thick] (2.5, 5.6) -- (5.0, 5.6) node[midway, above=2pt, fill=BG, font=\footnotesize\bfseries] {Уязвимый интервал ($\Delta t$)};

    % ==========================================
    % ОТРИСОВКА ПАКЕТОВ (\Delta t = 2.5)
    % ==========================================

    % --- СТАНЦИЯ 1 (Коллизия внутри слота k) ---
    % Пакет полностью совпадает с границами слота [5.0, 7.5]
    \fill[Red3!25] (5.0, 4.1) rectangle (7.5, 4.9);
    \draw[thick] (5.0, 4.1) rectangle (7.5, 4.9);
    \node[font=\footnotesize] at (6.25, 4.5) {Пакет $i$};
    
    \node[font=\footnotesize\bfseries, text=Red3, anchor=west] (col) at (7.8, 3.75) {Коллизия};
    % Стрелки "вилочкой" к конфликтующим пакетам
    \draw[->, thick, Red3] (col.west) -- (7.5, 4.4);
    \draw[->, thick, Red3] (col.west) -- (7.5, 3.1);

    % --- СТАНЦИЯ 2 (Целевой пакет внутри слота k) ---
    % Полностью перекрывается с Пакетом i
    \fill[Red3!25] (5.0, 2.6) rectangle (7.5, 3.4);
    \draw[thick] (5.0, 2.6) rectangle (7.5, 3.4);
    
    % Подпись вынесена за пределы пакета (сверху) для читаемости
    \node[font=\footnotesize, align=center] at (6.25, 3.0) {Целевой\\пакет};
    

    % --- СТАНЦИЯ 3 (Успешная передача в следующем слоте k+1) ---
    % Пакет строго в границах [7.5, 10.0]
    \fill[SpringGreen4!15] (7.5, 1.1) rectangle (10.0, 1.9);
    \draw[thick] (7.5, 1.1) rectangle (10.0, 1.9);
    \node[font=\footnotesize, align=center] at (8.75, 1.5) {Успешная\\передача};

    % --- СТАНЦИЯ 4 (Успешная передача в предыдущем слоте k-1) ---
    % Пакет строго в границах [2.5, 5.0]
    \fill[SpringGreen4!15] (2.5, -0.4) rectangle (5.0, 0.4);
    \draw[thick] (2.5, -0.4) rectangle (5.0, 0.4);
    \node[font=\footnotesize, align=center] at (3.75, 0) {Успешная\\передача};

\end{tikzpicture}
\end{center}

Уязвимый интервал для целевой передачи сужается до одного предыдущего слота $[(k-1)\Delta t, k\Delta t]$. Если на этом интервале длительностью $\tau = \Delta t$ ни один конкурирующий узел не сгенерирует заявку, в следующем слоте целевой пакет будет передан успешно.

Математическое ожидание числа событий фонового потока на интервале $\Delta t$ равно $\lambda \Delta t = G$. Вероятность того, что фоновых событий не произойдет, составляет $\mathbb{P}(k=0) = e^{-G}$.
Следовательно, нормированная пропускная способность равна $S = G e^{-G}$.

Точка экстремума определяется из уравнения:
$$\frac{dS}{dG} = e^{-G} - Ge^{-G} = e^{-G}(1 - G) = 0 \implies G = 1.$$
\end{sproof}

\begin{note}
При оптимальной нагрузке ($G=1$) вероятностное распределение состояний произвольного слота имеет вид:
\begin{itemize}
    \item Слот пуст (простой канала): $\P(0) = \frac{1^0}{0!}e^{-1} = e^{-1} \approx 0.368$.
    \item Успешная передача (ровно 1 пакет): $\P(1) = \frac{1^1}{1!}e^{-1} = e^{-1} \approx 0.368$.
    \item Коллизия ($k \ge 2$ пакетов): $\P(k \ge 2) = 1 - \P(0) - \P(1) = 1 - 2e^{-1} \approx 0.264$.
\end{itemize}
Данное распределение описывает безусловную статистику состояний канала и вычислено безотносительно к попытке передачи какого-либо выделенного пакета.
\end{note}

\begin{center}
\begin{tikzpicture}
\begin{axis}[
    width=0.85\textwidth, height=7.5cm,
    xlabel={Суммарная нагрузка $G$ (пакеты/слот)},
    ylabel={Пропускная способность $S$},
    xmin=0, xmax=3, ymin=0, ymax=0.45,
    grid=major,
    legend pos=north east,
    samples=150, thick
]
    % Асинхронная ALOHA
    \addplot[DeepSkyBlue4, domain=0:3] {x * exp(-2*x)};
    \addlegendentry{Асинхронная ALOHA}

    % Синхронная ALOHA
    \addplot[Red3, domain=0:3] {x * exp(-x)};
    \addlegendentry{Синхронная ALOHA}

    % Максимум асинхронной
    \draw[dashed, gray] (axis cs:0.5, 0) -- (axis cs:0.5, 0.1839) -- (axis cs:0, 0.1839);
    \node[circle, fill=DeepSkyBlue4, inner sep=1.5pt] at (axis cs:0.5, 0.1839) {};
    \node[above right, font=\footnotesize, DeepSkyBlue4] at (axis cs:0.5, 0.1839) {$(0.5; \frac{1}{2e})$};

    % Максимум синхронной
    \draw[dashed, gray] (axis cs:1, 0) -- (axis cs:1, 0.3678) -- (axis cs:0, 0.3678);
    \node[circle, fill=Red3, inner sep=1.5pt] at (axis cs:1, 0.3678) {};
    \node[above right, font=\footnotesize, Red3] at (axis cs:1, 0.3678) {$(1; \frac{1}{e})$};
\end{axis}
\end{tikzpicture}
\end{center}

\subsection{Синхронная ALOHA при конечном числе устройств}

В предыдущем разделе пропускная способность сети вычислялась в асимптотическом предположении о бесконечном числе терминалов ($n \to \infty$), генерирующих суммарный пуассоновский поток. Рассмотрим более реалистичную модель слотированной ALOHA с конечным числом узлов.

\begin{statement}[Пропускная способность системы с конечным числом терминалов]
Рассмотрим сеть, состоящую из $n$ независимых терминалов. Пусть $G_i \in (0, 1)$ --- вероятность начала передачи пакета $i$-м терминалом в произвольном слоте. Суммарная нагрузка на канал $G$ (математическое ожидание числа попыток передачи в одном слоте) равна $G = \dsum_{i=1}^{n} G_i$.

Индивидуальная пропускная способность $S_i$ (вероятность успешной передачи $i$-го терминала) и суммарная пропускная способность системы $S$ задаются соотношениями:
$$
S_i = \frac{G_i}{1 - G_i} \prod_{j=1}^{n} (1 - G_j), \qquad S = \sum_{i=1}^{n} S_i.
$$
\end{statement}

\begin{sproof}
Успешная передача пакета $i$-м терминалом в выделенном слоте происходит тогда и только тогда, когда $i$-й терминал осуществляет передачу, а все остальные $n-1$ терминалов «молчат» (не инициируют передачу). В силу независимости работы терминалов вероятность совместного наступления этих событий равна произведению их вероятностей:
$$
S_i = G_i \prod_{\substack{j=1 \\ j \neq i}}^{n} (1 - G_j).
$$
Для приведения выражения к симметричному виду относительно произведения по всем $j$, умножим и разделим правую часть на вероятность отсутствия передачи $i$-го терминала $(1 - G_i)$:
$$
S_i = \frac{G_i}{1 - G_i} \prod_{j=1}^{n} (1 - G_j).
$$
События, заключающиеся в успешной передаче пакета $i$-м или $k$-м терминалом в одном и том же слоте, являются несовместными (в противном случае возникает коллизия). Следовательно, суммарная пропускная способность сети вычисляется как сумма индивидуальных пропускных способностей:
$$
S = \sum_{i=1}^{n} S_i = \left( \sum_{i=1}^{n} \frac{G_i}{1 - G_i} \right) \prod_{j=1}^{n} (1 - G_j).
$$
\end{sproof}

\begin{conseq}
Предположим, что нагрузка распределена между терминалами равномерно. Тогда для любого $i \in \{1, \dots, n\}$ вероятность передачи составляет $G_i = \frac{G}{n}$. Подставим данное значение в формулу суммарной пропускной способности:
$$
S = n \cdot S_i = n \frac{\frac{G}{n}}{1 - \frac{G}{n}} \left(1 - \frac{G}{n}\right)^n = \frac{G}{1 - \frac{G}{n}} \left(1 - \frac{G}{n}\right)^n.
$$
Устремим число узлов к бесконечности ($n \to \infty$) при строго фиксированной суммарной нагрузке $G$. Используя классический предел $\lim\limits_{n \to \infty} \left(1 - \frac{x}{n}\right)^n = e^{-x}$ и учитывая, что знаменатель $\left(1 - \frac{G}{n}\right) \to 1$, вычисляем предел пропускной способности:
$$
\lim_{n \to \infty} S = G \lim_{n \to \infty} \left(1 - \frac{G}{n}\right)^n = G e^{-G}.
$$
Данный результат совпадает с формулой пропускной способности слотированной ALOHA с бесконечным числом устройств, полученной ранее на основе гипотезы о пуассоновском характере фонового трафика.
\end{conseq}
